% Provenance (never printed): SNT "Le Web" (client-server, HTTP, forms) carried into a project; NSI premiere
% "Interactions entre l'homme et la machine sur le Web" (forms, GET/POST, server side). See sources/book2/BRIEFS.md.
\chapter{Project: A Small Web Application}\label{ch:g10:project-web-application}

The chapter on \omterm{def:g10:http-clients-servers:http}{HTTP} ended with a server that knew five addresses by heart. A real site is not written that way:
its pages are built from data at the moment they are asked for, its forms change that data, and it is tested as
carefully as any other program. This project builds such a site for the Brenmoor library, from a one-paragraph
specification to a tested program, in five milestones.

\begin{recall}
A request has a method, a path and headers; the answer has a \omterm{def:g10:http-clients-servers:status}{status code}, headers and a body. \omterm{def:g10:http-clients-servers:methods}{GET} puts its
parameters in the \omterm{def:g10:http-clients-servers:url}{query string}, \omterm{def:g10:http-clients-servers:methods}{POST} in the body; a \omterm{def:g10:http-clients-servers:cookie}{cookie} is a value the browser sends back to the server that
set it (the chapter on \omterm{def:g10:http-clients-servers:http}{HTTP}). A join matches the rows of two tables on a common column (the chapter on tables). A
test compares what a function does with what it should do (the chapter on testing). An HTML page is made of
elements written as tags, and a style sheet decides how they look (Book 1, grade 9).
\end{recall}

\section{The specification}\label{sec:g10:project-web-application:spec}

\begin{quote}
\itshape A catalogue of the Brenmoor library, served on \code{127.0.0.1} at port $42110$, built from the town's
tables of books, loans and members. It has a home page listing every book; a page for each book, with its loans; a
search form that finds books by a word of their title; a form on each book's page to reserve it, given a member
number, which records the reservation in memory; and a clear error page for any other address. Every address is
tested. Nothing leaves the machine.
\end{quote}

The program is split in three layers. A small \emph{server} part speaks \omterm{def:g10:http-clients-servers:http}{HTTP}: it turns each request into a
dictionary and sends back what the application returns. The \emph{application} decides, for each path, which
function builds the answer. The \emph{data} part reads the three tables once, at the start.

\begin{omfigure}
\begin{tikzpicture}[font=\footnotesize, b/.style={draw, rounded corners=2pt, minimum height=10mm, align=center,
    minimum width=24mm}]
  \node[b, fill=ocBlue!10] (br) at (0,0) {browser\\or test};
  \node[b, fill=omCode!10] (sv) at (4.0,0) {\code{server.py}\\request $\to$ dict};
  \node[b, fill=omDef!10] (ap) at (8.2,0) {\code{dispatch}\\route $\to$ function};
  \node[b, fill=omMeth!10] (da) at (11.8,0.7) {\code{data.py}\\three tables};
  \node[b, fill=omProp!10] (re) at (11.8,-0.7) {reservations\\(in memory)};
  \draw[ptr, <->] (br) -- node[above] {HTTP} (sv);
  \node[font=\scriptsize, below] at (2.0,-0.55) {port 42110};
  \draw[ptr, ->] ([yshift=2mm]sv.east) -- node[above] {dict} ([yshift=2mm]ap.west);
  \draw[ptr, ->] ([yshift=-2mm]ap.west) -- node[below] {answer} ([yshift=-2mm]sv.east);
  \draw[ptr] (ap.east) -- (da.west);
  \draw[ptr, <->] (ap.east) -- (re.west);
  \node[font=\scriptsize, black!60] at (11.8,1.6) {\code{books.csv} \code{loans.csv} \code{members.csv}};
\end{tikzpicture}

{\small The three layers. Only the server part knows about \omterm{def:g10:http-clients-servers:http}{HTTP}; the application receives a dictionary and returns
a status, a type, a text and some headers, so it can be tested without any network at all.}
\end{omfigure}

The server part is given: \code{server.py} builds a dictionary with the method, the path (with its \texttt{\%}
escapes decoded), the query, the form and the \omterm{def:g10:http-clients-servers:cookie}{cookies}, calls \code{app.dispatch}, and writes the answer. A function
\code{start(app, port)} runs it in the background, and \code{client.py} sends one request and shows the answer
without following redirects.

\omcode{code/grade-10/11-project-web-application/listings/server.py}{12}{34}{The heart of the given server: a
request becomes a dictionary.}\label{lst:g10:project-web-application:server}

\section{Milestone 1: routes and a dispatcher}\label{sec:g10:project-web-application:m1}

\begin{definition}[Route, dispatcher]\label{def:g10:project-web-application:route}
In a web application, a \emph{route}\index{route (web application)} links a family of paths to the function that
answers them, such as every path \code{/book/<number>} to the function \code{book}. The
\emph{dispatcher}\index{dispatcher} is the function that receives each request, finds its route, and calls that
function with the parts of the path it needs.
\end{definition}

The first milestone answers in plain text: the point is to get the \omterm{def:g10:project-web-application:route}{routes} right before any page is built. The \omterm{def:g10:project-web-application:route}{route}
table is a dictionary from the first part of the path to a pair: the function, and the number of further parts it
takes. The key \code{None} holds the function for every path that matches no \omterm{def:g10:project-web-application:route}{route}.

\omcode{code/grade-10/11-project-web-application/listings/app_m1.py}{1}{34}{Milestone 1: a \omterm{def:g10:project-web-application:route}{route} table and a
\omterm{def:g10:project-web-application:route}{dispatcher}.}\label{lst:g10:project-web-application:m1}

\omcode{code/grade-10/11-project-web-application/listings/run_m1.py}{1}{9}{Start, request six paths,
stop.}\label{lst:g10:project-web-application:run-m1}

\omoutput{code/grade-10/11-project-web-application/outputs/run_m1.txt}{Every path gets an answer: a page, or a
404 that says what was asked.}\label{out:g10:project-web-application:run-m1}

\begin{omfigure}
\begin{tikzpicture}[font=\footnotesize, p/.style={draw, fill=black!4, minimum width=34mm, minimum height=6mm,
    anchor=west, font=\ttfamily\footnotesize}, f/.style={draw, fill=omDef!10, rounded corners=2pt,
    minimum width=22mm, minimum height=6mm, anchor=west}]
  \foreach \y/\pth/\fun/\m in {0/{/}/home/1, -0.75/{/book/<n>}/book/1, -1.5/{/style.css}/style/2,
      -2.25/{/search?q=\dots}/search/3, -3/{/reserve \ (POST)}/reserve/3, -3.75/{/api/book/<n>}/api/5,
      -4.5/{anything else}/not\_found/1} {
    \node[p] (a) at (0,\y) {\pth};
    \node[f, font=\ttfamily\footnotesize] (b) at (5,\y) {\fun};
    \draw[ptr] (a.east) -- (b.west);
    \node[anchor=west, font=\scriptsize, black!60] at (7.6,\y) {milestone \m};
  }
\end{tikzpicture}

{\small The \omterm{def:g10:project-web-application:route}{route} map of the finished application, with the milestone that adds each \omterm{def:g10:project-web-application:route}{route}. Milestones 3 and 5
also replace the functions for \code{/book/<n>}, \code{/} and \code{/search} with richer ones.}
\end{omfigure}

\section{Milestone 2: pages built from the tables}\label{sec:g10:project-web-application:m2}

The pages are now HTML, built by a function \code{page(title, body)} that adds the parts every page shares. A book's
page joins two tables: the book's row in \code{BOOKS}, and its loans in \code{LOANS}, a dictionary from book number
to loans that \code{data.py} builds once at the start.

A page is text, and some of that text comes from outside: a title in a table, a word typed in a form, a path. If a
path contained \code{<script>}, a page that printed it as it is would contain a script, and the browser would run
it. Text that comes from outside must never be read as HTML.

\begin{definition}[Escaping]\label{def:g10:project-web-application:escaping}
\emph{Escaping}\index{escaping (HTML)} a text for HTML replaces each character that has a meaning in HTML by an
entity that only displays it: \code{<} becomes \code{&lt;}, \code{>} becomes \code{&gt;}, \code{&} becomes
\code{&amp;}, and the quotes become \code{&quot;} and \texttt{\&\#x27;}. Python's \code{html.escape} does it.
\end{definition}
% ledger: g10:html-escape

\omcode{code/grade-10/11-project-web-application/listings/app_m2.py}{1}{36}{Milestone 2: HTML pages, every value
from outside escaped.}\label{lst:g10:project-web-application:m2}

The rest of the file adds \code{not_found}, which escapes the path it shows, a \omterm{def:g10:project-web-application:route}{route} for the style sheet, and a
\code{dispatch} that calls the one of milestone~1 with the new \omterm{def:g10:project-web-application:route}{route} table: the \omterm{def:g10:project-web-application:route}{dispatcher} itself does not change.

\omcode{code/grade-10/11-project-web-application/listings/run_m2.py}{1}{11}{The home page, a book's page, and an
attack on the error page.}\label{lst:g10:project-web-application:run-m2}

\omoutput{code/grade-10/11-project-web-application/outputs/run_m2.txt}{The script in the path comes back as text:
the browser shows it, and runs nothing.}\label{out:g10:project-web-application:run-m2}

\omcode{code/grade-10/11-project-web-application/listings/style.css}{1}{5}{The style sheet, served at
\code{/style.css}.}\label{lst:g10:project-web-application:style}

\begin{omfigure}
\begin{tikzpicture}
  \omwindow{8}{4.6}{The Silent Kingdom}
  \pic at (0,0) {window};
  \node[anchor=north west, font=\sffamily\bfseries\large, text=ocBlue] at (0.3,3.95) {The Silent Kingdom};
  \node[anchor=west, font=\sffamily\small] at (0.3,3.0) {Ines Yewwell, 1950, 286 pages};
  \foreach \i/\o/\b in {0/2025-02-09/2025-02-24, 1/2025-03-20/2025-03-30, 2/{\dots}/{\dots},
      3/2025-07-16/not yet} {
    \draw[black!50] (0.3,2.5-0.35*\i) rectangle (1.9,2.15-0.35*\i);
    \draw[black!50] (1.9,2.5-0.35*\i) rectangle (3.5,2.15-0.35*\i);
    \node[font=\sffamily\scriptsize] at (1.1,2.325-0.35*\i) {\o};
    \node[font=\sffamily\scriptsize] at (2.7,2.325-0.35*\i) {\b};
  }
  \node[anchor=west, font=\sffamily\small] at (0.3,0.85) {Reserved by: 12};
  \node[anchor=west, font=\sffamily\small] at (0.3,0.35) {Member number:};
  \draw[black!60, fill=white] (3.0,0.15) rectangle (4.8,0.55);
  \node[draw=black!60, fill=black!8, rounded corners=1pt, font=\sffamily\small] at (5.55,0.35) {Reserve};
\end{tikzpicture}

{\small A book's page as a browser shows it, from the HTML of milestone~3: the loans table (rows left out here), the
reservations, and the form that sends \code{book} and \code{member} to \code{/reserve} with \omterm{def:g10:http-clients-servers:methods}{POST}.}
\end{omfigure}

\section{Milestone 3: forms}\label{sec:g10:project-web-application:m3}

The search form uses \omterm{def:g10:http-clients-servers:methods}{GET}: searching changes nothing on the server, and its address can be bookmarked. The
reservation form uses \omterm{def:g10:http-clients-servers:methods}{POST}: it changes the server's data. After a successful \omterm{def:g10:http-clients-servers:methods}{POST}, the server does not send a
page; it answers \code{303 See Other} with a \code{Location} header, and the browser then asks for the book's page
with \omterm{def:g10:http-clients-servers:methods}{GET}. This is called \emph{redirect after \omterm{def:g10:http-clients-servers:methods}{POST}}: if the reader reloads the page, the browser repeats the \omterm{def:g10:http-clients-servers:methods}{GET},
not the reservation.
% ledger: g10:http-9110-classes

\omcode{code/grade-10/11-project-web-application/listings/app_m3.py}{13}{43}{Milestone 3: a search, a book page
with a form, and the reservation.}\label{lst:g10:project-web-application:m3}

Every value of the form is checked before it is used: a book number that is not a number, or not in the catalogue,
and a member number that is not a member, are answered with \code{400 Bad Request}. A program that trusts its forms
can be broken by anyone who types into them, or who sends a request without any form at all.

\omcode{code/grade-10/11-project-web-application/listings/run_m3.py}{1}{11}{Two searches, two reservations, and
the book's page.}\label{lst:g10:project-web-application:run-m3}

\omoutput{code/grade-10/11-project-web-application/outputs/run_m3.txt}{The second reservation names a member who
does not exist, and is refused.}\label{out:g10:project-web-application:run-m3}

\begin{omfigure}
\begin{tikzpicture}[font=\footnotesize, >={Stealth[length=2mm]}]
  \node[font=\bfseries\footnotesize] at (0,0) {browser};
  \node[font=\bfseries\footnotesize] at (7,0) {server};
  \draw[black!40, dashed] (0,-0.3) -- (0,-5.6);
  \draw[black!40, dashed] (7,-0.3) -- (7,-5.6);
  \foreach \y/\l/\d in {-0.7/{GET /book/7}/1, -1.3/{200: the page, with its form}/0,
      -2.1/{POST /reserve \ book=7\&member=12}/1, -2.7/{303 See Other, Location: /book/7}/0,
      -3.5/{GET /book/7}/1, -4.1/{200: Reserved by: 12}/0, -4.9/{reload: GET /book/7 again}/1} {
    \ifnum\d=1
      \draw[->, thick, ocBlue] (0,\y) -- node[above, text=black] {\l} (7,\y);
    \else
      \draw[->, thick, omProp] (7,\y) -- node[above, text=black] {\l} (0,\y);
    \fi
  }
  \node[anchor=west, font=\scriptsize, black!60] at (7.2,-4.9) {no second reservation};
\end{tikzpicture}

{\small Redirect after \omterm{def:g10:http-clients-servers:methods}{POST}. The \omterm{def:g10:http-clients-servers:methods}{POST} is answered by a 303, and the browser itself asks for the page with \omterm{def:g10:http-clients-servers:methods}{GET};
reloading that page repeats only the \omterm{def:g10:http-clients-servers:methods}{GET}.}
\end{omfigure}

\section{Milestone 4: tests against the real server}\label{sec:g10:project-web-application:m4}

Each milestone has been tried by hand. A \omterm{def:g10:specifications-and-tests:test-suite}{test suite} tries every \omterm{def:g10:project-web-application:route}{route} at once, every time the program changes. It
starts the real server in a process of its own, sends \omterm{def:g10:http-clients-servers:http}{HTTP} requests with \code{urllib} like any client, and checks
each status and one line of each page. \code{urlopen} follows the 303 by itself, and turns a 4xx or 5xx answer into
an exception, whose code the helper \code{get} returns.

\omcode{code/grade-10/11-project-web-application/listings/app_tests.py}{1}{40}{Milestone 4: the server under test,
started once for all the tests.}\label{lst:g10:project-web-application:tests-setup}

\omcode{code/grade-10/11-project-web-application/listings/app_tests.py}{42}{75}{The tests: one per \omterm{def:g10:project-web-application:route}{route}, and the
errors.}\label{lst:g10:project-web-application:tests}

\omoutput{code/grade-10/11-project-web-application/outputs/run_tests.txt}{The command is
\code{python -m unittest -v app_tests}; the line that gives the duration is left out.}\label{out:g10:project-web-application:tests}

\begin{inpractice}[Two kinds of tests]
The tests of this milestone go through the network, the server and the application: they are slow, a fraction of
a second each, but they would catch a mistake in any layer. Because the application only receives a dictionary,
it can also be tested by calling \code{dispatch} directly, with no server at all; the chapter's own tests do both.
\end{inpractice}

\section{Milestone 5: a cookie and a JSON route}\label{sec:g10:project-web-application:m5}

Two improvements finish the project. The search answer sets a \omterm{def:g10:http-clients-servers:cookie}{cookie} \code{last}, which the browser sends back with
every later request; the home page then offers the last search as a link. And a \omterm{def:g10:project-web-application:route}{route} \code{/api/book/<n>} answers
in \omterm{def:g10:data-tables:json}{JSON} instead of HTML, for another program rather than a person: the same data, in the format of the chapter on
tables.

\omcode{code/grade-10/11-project-web-application/listings/app_m5.py}{13}{45}{Milestone 5: wrapping the search
and the home page, and a \omterm{def:g10:project-web-application:route}{route} for programs.}\label{lst:g10:project-web-application:m5}

\omcode{code/grade-10/11-project-web-application/listings/run_m5.py}{1}{14}{A search, the home page with the
\omterm{def:g10:http-clients-servers:cookie}{cookie}, and the \omterm{def:g10:data-tables:json}{JSON} \omterm{def:g10:project-web-application:route}{route}.}\label{lst:g10:project-web-application:run-m5}

\omoutput{code/grade-10/11-project-web-application/outputs/run_m5.txt}{The \omterm{def:g10:http-clients-servers:cookie}{cookie}'s value is percent-encoded, like
a \omterm{def:g10:http-clients-servers:url}{query string}; the \omterm{def:g10:data-tables:json}{JSON} names its fields.}\label{out:g10:project-web-application:run-m5}

The application meets its specification: every \omterm{def:g10:project-web-application:route}{route} answers, every value from outside is checked or escaped,
and the \omterm{def:g10:specifications-and-tests:test-suite}{test suite} proves it each time it runs. What it lacks is memory: the reservations live in a dictionary and
vanish when the server stops. Keeping them in a database, which several servers can share safely, is the subject of
a later chapter.

\section{Exercises}\label{sec:g10:project-web-application:exercises}

\begin{exercise}[$\star$]\label{exo:g10:project-web-application:1}
For each path, say which function of milestone~1 answers, and with which arguments: \code{/}, \code{/book/12},
\code{/book/12/}, \code{/book/12/loans}, \code{/books/12}.
\end{exercise}

\begin{exercise}[$\star$]\label{exo:g10:project-web-application:2}
What does the page of milestone~2 contain, exactly, for the path \code{/a&b<i>}? What would a browser display?
\end{exercise}

\begin{exercise}[$\star$]\label{exo:g10:project-web-application:3}
Add a \omterm{def:g10:project-web-application:route}{route} \code{/about} with a page of one sentence about the catalogue.
\end{exercise}

\begin{exercise}[$\star$]\label{exo:g10:project-web-application:4}
How many \omterm{def:g10:http-clients-servers:http}{HTTP} requests does the \omterm{def:g10:specifications-and-tests:test-suite}{test suite} of milestone~4 send in all? Count the redirection.
\end{exercise}

\begin{exercise}[$\star$]\label{exo:g10:project-web-application:5}
A student forgets the word \code{return} on the last line of \code{book}. What does \code{dispatch} give back for
\code{/book/7}, and what happens in the server?
\end{exercise}

\begin{exercise}[$\star\star$]\label{exo:g10:project-web-application:6}
Add a \omterm{def:g10:project-web-application:route}{route} \code{/author/<name>} that lists the books of an author, with a 404 if there are none. How many books
does Zoe Stoneton have?
\end{exercise}

\begin{exercise}[$\star\star$]\label{exo:g10:project-web-application:7}
In a copy of \code{reserve}, the values are read from \code{request["query"]} instead of
\code{request["form"]}. What does every reservation sent by the form now get, and why? Which test catches it?
\end{exercise}

\begin{exercise}[$\star\star$]\label{exo:g10:project-web-application:8}
Let the home page take a parameter \code{sort}, which may be \code{title}, \code{author} or \code{year}; any other
value gets a 400.
\end{exercise}

\begin{exercise}[$\star\star$]\label{exo:g10:project-web-application:9}
With $n$ books in the catalogue, how many different addresses answer 200 with \omterm{def:g10:http-clients-servers:methods}{GET}, if searches are not counted? How
many for the Brenmoor catalogue?
\end{exercise}

\begin{exercise}[$\star\star$]\label{exo:g10:project-web-application:10}
A test file starts the server with \code{subprocess.Popen} at its top, runs its tests, and stops the server on its
last line. One test fails. What happens to the server, and to the next run of the tests? Fix it.
\end{exercise}

\begin{exercise}[$\star\star$]\label{exo:g10:project-web-application:11}
The specification says that \code{/reserve} accepts only \omterm{def:g10:http-clients-servers:methods}{POST}, and the other \omterm{def:g10:project-web-application:route}{routes} only \omterm{def:g10:http-clients-servers:methods}{GET}. Make the \omterm{def:g10:project-web-application:route}{dispatcher}
answer \code{405 Method Not Allowed} otherwise, with the header that this status requires.
\end{exercise}

\begin{exercise}[$\star\star\star$]\label{exo:g10:project-web-application:12}
Show the home page twenty books at a time: \code{/?page=1} to the last page, 404 outside. Test the first and last
pages, the pages just outside, and that all the pages together show every book once.
\end{exercise}

\begin{exercise}[$\star\star\star$]\label{exo:g10:project-web-application:13}
A reader reserves a book, then reloads the page three times. Count the reservations recorded if the server answers
the \omterm{def:g10:http-clients-servers:methods}{POST} with the book's page directly, and if it answers with a 303. Why does a browser warn before repeating a
\omterm{def:g10:http-clients-servers:methods}{POST}?
\end{exercise}

\begin{exercise}[$\star\star\star$]\label{exo:g10:project-web-application:14}
A member may hold at most three reservations. Refuse a fourth with status \code{409 Conflict}, and test it.
\end{exercise}

\begin{exercise}[$\star\star\star$]\label{exo:g10:project-web-application:15}
The server answers several requests at the same time, one thread each. Member $5$ holds two reservations and sends
two more at the same moment. Show an order of the steps of the limit check of the previous exercise in which both
are accepted. How can the program prevent it?
\end{exercise}
